\section{Related Work}
\label{sec:related}

\subsection{Affected-Version Identification}

Affected-version identification takes a disclosed vulnerability and predicts the set of released versions in which the vulnerable logic remains present. Chen et al. formalize this task over a release universe and evaluate it at two granularities~\cite{howfar2025}. Vulnerability-level evaluation requires an exact match of the complete affected-version set. Version-level precision, recall, and F1 instead measure individual CVE-version decisions and therefore capture partial correctness. This distinction separates affected-version identification from vulnerability-inducing commit (VIC) localization and from advisory range normalization.

Direct methods use different forms of repository evidence. V-SZZ traces patch-related lines to candidate VICs, detects duplicated changes across branches, and maps inducing and fixing commits to version tags~\cite{vszz2022}. The patch-and-developer-log approach of He et al. organizes releases as a version tree and reasons about first and last affected versions, branch boundaries, partial matches, and repatching~\cite{tdsc2024}. Vercation combines program slicing, LLM-based vulnerable-statement selection, and structural clone comparison before deriving affected versions between a VIC and a fixing commit~\cite{vercation2026}. A recent anonymous manuscript, \emph{CaVulner}, also uses CVE/CWE context, LLM-assisted root-cause localization, blame, and duplicate-change analysis.\footnote{Citation gap: the available CaVulner manuscript is anonymous and has no stable publication metadata.}

These methods establish that patch semantics, commit history, duplicated fixes, and release structure all matter. However, a VIC, a fixing commit, or a version-range endpoint does not by itself represent partial fixes, equivalent fixes, reversions, and multiple branch-specific state changes. VulnVersion therefore targets complete released-version set prediction through history event reconstruction, boundary judgment, Git-DAG vulnerability-state propagation, and release tag projection.

\subsection{SZZ and Agent-Assisted History Reconstruction}

SZZ research improves the evidence used to locate bug-inducing commits (BICs) or VICs. Early approaches derive candidates from lines changed by a fixing commit. VCCFinder uses blame-derived labels together with code and repository features to prioritize suspicious commits, while Lifetime uses weighted blame evidence to estimate vulnerability lifetimes~\cite{vccfinder2015,lifetime2021}. Semantic variants refine which statements should be traced. SEM-SZZ compares local program states and data-flow evidence; TC-SZZ performs iterative history tracing and studies cases that direct blame cannot recover; and differential patch-pattern analysis extracts vulnerability-critical statement sequences before matching earlier commits~\cite{semszz2024,tcszz2024,vicpatterns2026}. These methods improve candidate evidence, but their main output remains a commit or a lifetime boundary.

LLM-assisted methods replace fixed heuristics with semantic assessment. LLM4SZZ ranks buggy statements and evaluates candidate commits with expanded context~\cite{llm4szz2025}. AgentSZZ adds task-specific repository tools, domain knowledge, and iterative exploration for ghost and cross-file cases~\cite{agentszz2026}. Other agentic workflows search file histories through binary or direct candidate selection and derive short patterns from the fixing commit for searching candidate changes~\cite{howwhyagents2026}. These approaches broaden repository exploration, but they are evaluated on BIC identification rather than complete affected-version set prediction.

MAS-SZZ is especially close to the evidence layer of VulnVersion. It summarizes a root cause with traceable evidence, infers patch-hunk intent, selects vulnerability-related anchor statements, and uses an agent to backtrack until it finds a VIC~\cite{masszz2026}. VulnVersion does not claim that evidence-grounded root-cause analysis or semantic anchor selection is unique. Its distinction begins after anchor selection: selected statements act as evidence-constrained pre-fix anchors for reconstructing history event chains and judging branch-specific introduction boundaries, rather than serving only as paths to a final VIC.

Beyond Blame is also closely related because it reframes BIC identification as agent-guided search over a temporal knowledge graph. Its graph expands candidates from blame commits and the fixing commit through temporal and structural commit relations~\cite{beyondblame2026}. VulnVersion shares the view that a single blame result is insufficient. The task and graph roles differ, however. Beyond Blame searches for BICs; VulnVersion reconstructs introduction, transformation, fix, and reversion events, then propagates vulnerability state through the Git DAG and projects that state to official release tags.

\subsection{Matching-Based and Recurring Vulnerability Detection}

Matching-based systems provide direct evidence that vulnerable or patched code occurs in a target codebase. ReDeBug uses normalized token windows and Bloom filters for scalable unpatched-clone search, while VUDDY indexes normalized function fingerprints~\cite{redebug2012,vuddy2017}. These designs provide efficient and precise evidence for exact or abstracted clones. Their function or token granularity, however, does not encode the full repository context in which a matched fragment is security-relevant.

Modification-aware systems recover more varied forms of reuse. MOVERY distinguishes vulnerability signatures from patch signatures, V1SCAN combines version-based filtering with code classification, and VULTURE combines version evidence with patch- and chunk-level analysis~\cite{movery2022,v1scan2023,vulture2025}. FIRE further uses staged clone filters followed by taint-based verification~\cite{fire2024}. These methods should not be reduced to simple textual matching: they contribute normalized signatures, structural evidence, patch-state evidence, code classification, and taint paths.

Their remaining limitation is a task boundary. A detector centered on deleted pre-fix text has little direct evidence for an add-only guard. Exact and near-exact signatures can also lose matches after semantic transformation, code movement, or branch divergence, although abstraction and modification-aware signatures mitigate some of these cases. Conversely, a code match does not establish that trigger conditions, guards, and reachable sinks are unchanged. Matching evidence is therefore valuable for candidate generation and boundary judgment, but it cannot alone reconstruct branch-specific affected versions across a Git DAG.

\subsection{Security Knowledge Graphs and Attacker-Condition Reasoning}

Security knowledge graphs commonly organize threat intelligence rather than repository evolution. ATT\&CK-to-CVE graphs integrate CVE, CWE, CAPEC, and ATT\&CK entities to support cross-source and multi-hop analysis~\cite{attacktocve2025}. NEXUS maps CVE descriptions to ATT\&CK techniques and supports analyst feedback for later mapping decisions~\cite{nexus2026}. D3FEND models defensive techniques, artifacts, and their relations to offensive techniques, while KRYSTAL links audit provenance with background knowledge for attack detection and scenario reconstruction~\cite{d3fend2021,krystal2022}. These graphs answer questions about attack techniques, defenses, threat intelligence, or observed system events. They do not reconstruct affected release states from a fixing commit.

Beyond Blame uses a different kind of graph: a temporal graph of commits for BIC search~\cite{beyondblame2026}. This graph preserves temporal and structural repository relations, but it is not an attacker-condition graph. The distinction matters because an affected-version decision may depend on attacker-controlled input, an exploit precondition, a reachable sink, a missing guard, or an unsafe state even when similar code is present.

For VulnVersion, the attacker perspective defines which security conditions should be preserved as structured evidence during history reconstruction. It does not mean generating or executing an exploit for every release tag. The intended graph role is to connect attacker conditions and root-cause evidence to pre-fix anchors, history events, and boundary events. We do not claim in this paper that an attacker-condition knowledge graph has already improved affected-version results; that claim requires a dedicated evaluation.

\subsection{Evidence-Gated Continual Adaptation}

Recent agent research studies long-term memory, procedural memory, and reusable skills. Surveys organize the roles and risks of agent memory, while procedural-memory systems extract reusable action knowledge from prior trajectories~\cite{agentmemorysurvey2025,memps2025}. Self-evolving agent work further learns, verifies, or distills reusable skills from experience~\cite{memskill2026,autoskill2026,coevo2026,trace2skill2026}. These directions motivate knowledge reuse across repeated software-analysis tasks.

Affected-version identification imposes a stricter evidence contract on such reuse. A prior decision should be reusable only when it is linked to verified repository evidence and records its provenance. Useful memory includes inspected failure patterns and repository-, CWE-, or patch-pattern priors. Ground-truth affected-version labels must remain separate from this memory, and contradictory evidence must be able to invalidate or downgrade an earlier observation. This evidence-gated continual adaptation differs from accepting free-form model memory as fact.

VulnVersion treats continual adaptation as a design boundary, not as an empirically validated contribution. Reused evidence may guide pre-fix anchor selection, history event ranking, and failure diagnosis, but it cannot replace boundary-event judgment or release tag projection. A dedicated ablation is required before any effectiveness claim is made.
% TODO(citation gap): The local Paper/reference corpus does not yet contain full-text packages for the cited agent-memory and skill-evolution papers. Before submission, verify detailed mechanism-level comparisons against their original texts.
